
\section{1. Präsenzübung 2024}
\subsection{Primitive Datentypen}
Die Menge aller valider Zeichenketten für einen primitiven Datentyp ist also eine formale Sprache. Insbesondere sind diese formalen Sprachen regulär, wie wir im folgenden beweisen werden.
\begin{itemize}
\item \textbf{Reguläre Ausdrücke} sind spezielle Formeln, die reguläre Sprachen beschreiben könnnen. \\ \newline
\item \textbf{Formale Sprachen} sind Sprachen, die durch Regeln formalisiert sind. \\ \newline
\item \textbf{Reguläre Sprachen} sind formale Sprachen, die durch endliche Automaten erkannt werden können.  \\ \newline
\item \textbf{Rechtslineare Grammatiken} haben Ersetzungsregeln der Form $A ::= \overline{w}B$ oder $A ::= \overline{w}$, wobei $\overline{w}$ ein Terminalsymbol oder meherer Terminalsymbole (ein Wort) sein kann.
\end{itemize}

\subsubsection*{Entwickeln Sie einen regulären Ausdruck, der die Sprache aller möglichen Boolschen Konstanten definiert.}

\[
 \Sigma = \kl{\un{a},\un{e}, \un{f}, \un{l}, \un{r}, \un{s},\un{t},\un{u}}
\]
\[
regex = \un{true} \mid \un{false}
\]

\subsubsection*{Entwickeln Sie einen regulären Ausdruck, der die Sprache aller möglichen Integer-Konstanten definiert.}

\[
\Sigma = \kl{\un{0}, ..., \un{9}}
\]

\[
regex = (\un{0} \mid .... \mid \un{9})+
\]

\subsubsection*{Entwickeln Sie einen regulären Ausdruck, der die Sprache aller möglichen String-Konstanten über dem Alphabet $\kl{\un{a}...\un{z}}$ definiert}

\[
\Sigma = \kl{\un{a}...\un{z}}
\]

\[
regex = (\un{a} \mid .... \mid \un{z})+
\]
\subsubsection*{Entwickeln Sie eine rechtslineare Grammatik, die die Sprache der Gleitkommazahlen Zahlen in E-Notation mit beliebig vielen (aber endlich vielen) Dezimalstellen erzeugt}

\begin{align*}
    G = \{ \quad \\
    & \Sigma = \kl{\un{-}, \un{.}, \un{E}, \un{0}, ..., \un{9}}, \\
    & \Phi = \kl{A, B, C,D,E}, \\
    & S = A,  \\
    & R = \kl{A :: = \un{-}B \mid \un{ \epsilon} B, \quad B ::= \un{0.}C \mid ... \mid \un{0.}C, \quad C::= \un{0}C \mid \un{9}C \mid  \un{0}D \mid \un{9}D, \quad D ::= \un{E-}E \mid \un{E}E, \quad E::= \un{0} \mid ... \mid \un{9}} \\
    \}
\end{align*}

\subsubsection*{Leiten Sie mit der Grammatik das Beispielwort $\un{-7.12E-1}$ ab.}

\[
\un{-}B \rightarrow \un{-7.}C \rightarrow \un{-7.1}C \rightarrow \un{-7.12}D \rightarrow \un{-7.12E-}E \rightarrow \un{-7.12E-1} \]
